\chapter{Catálogo de requisitos}
\label{ap:requisitos}

Este apéndice recoge el catálogo completo de requisitos de la segunda fase.
El capítulo de especificación resume los objetivos principales; aquí se deja
la lista completa para poder consultar cada requisito, su prioridad MoSCoW y
la evidencia con la que se valida en las pruebas.

\section{Requisitos funcionales}

\begin{small}
\begin{itemize}
\item \textbf{RF-01: Carga dinámica de modelos.} Cargar dinámicamente los modelos
de decisión generados por la primera fase, descubriéndolos sin intervención
manual entre sesiones. \textit{Prioridad: Must}.
\item \textbf{RF-02: Generación de entornos.} Generar especificaciones de entorno
de simulación a partir de descripciones en lenguaje natural.
\textit{Prioridad: Must}.
\item \textbf{RF-03: Ejecución de simulaciones.} Ejecutar simulaciones con uno o
varios modelos en el mismo entorno, soportando comparativas multi-modelo.
\textit{Prioridad: Must}.
\item \textbf{RF-04: Registro estructurado.} Registrar de forma estructurada los
eventos, episodios y trayectorias de decisión producidos durante cada
simulación. \textit{Prioridad: Must}.
\item \textbf{RF-05: Identificación de patrones.} Identificar patrones de
comportamiento en los resultados, incluyendo pasos concretos o eventos
críticos, y contrastarlos con el Knowledge Backbone.
\textit{Prioridad: Must}.
\item \textbf{RF-06: Visualizaciones analíticas.} Generar, como parte del
análisis, gráficas y visualizaciones sobre trayectorias y métricas para
incorporarlas al informe final. \textit{Prioridad: Should}.
\item \textbf{RF-07: Generación de informes.} Generar informes en PDF con los
resultados, la evidencia visual y las referencias a los paradigmas
analizados. \textit{Prioridad: Must}.
\item \textbf{RF-08: Persistencia de observaciones.} Persistir las observaciones
de simulación en el Knowledge Backbone como memorias reutilizables en
ejecuciones posteriores. \textit{Prioridad: Must}.
\item \textbf{RF-09: Historial de experimentos.} Conservar el historial completo de
experimentos y permitir su consulta, filtrado y comparativa, también mediante
preguntas en lenguaje natural resolubles con SQL. \textit{Prioridad: Should}.
\item \textbf{RF-10: Interacción conversacional.} Permitir la interacción con el
sistema mediante un chat conversacional gestionado por el Orchestrator, con
continuidad de la conversación dentro de la sesión. \textit{Prioridad: Must}.
\item \textbf{RF-11: Recomendaciones del Orchestrator.} Permitir al Orchestrator
recomendar al usuario qué hacer a continuación, qué modelos encajan y qué
herramientas tiene disponibles. \textit{Prioridad: Could}.
\item \textbf{RF-12: Replay de simulación.} Ofrecer una visualización del entorno
de simulación con \textit{replay} paso a paso y paneles de actividad por
agente. \textit{Prioridad: Should}.
\item \textbf{RF-13: Inspección del Knowledge Backbone.} Exponer vistas de
inspección del Knowledge Backbone (grafo, memorias y procedencia) accesibles
desde la propia interfaz web. \textit{Prioridad: Could}.
\end{itemize}
\end{small}

\vspace{0.5\baselineskip}

\section{Requisitos no funcionales}

\begin{small}
\begin{itemize}
\item \textbf{RNF-01: Modularidad y desacoplamiento.} La integración con la primera
fase debe realizarse mediante un contrato mínimo (\textit{duck typing}),
sin imponer dependencias de paquetes ni jerarquías de clases. Mientras
RF-01 fija el \emph{qué} (cargar los modelos sin intervención manual),
este requisito fija el \emph{cómo} (sin acoplar la evolución de las dos
fases).
\item \textbf{RNF-02: Extensibilidad de modelos.} La incorporación de nuevos
modelos generados por la primera fase no debe requerir modificaciones en
el código de la segunda fase.
\item \textbf{RNF-03: Trazabilidad.} Cada ejecución debe conservar su
configuración, sus eventos y sus artefactos finales con un identificador
único, de modo que sea posible reproducirla o auditarla.
\item \textbf{RNF-04: Reactividad.} El cliente web debe recibir actualizaciones del
estado de la simulación y de la actividad de cada agente en tiempo real,
sin \textit{polling}, mediante WebSocket.
\item \textbf{RNF-05: Degradación controlada.} Si la capa semántica opcional
(\textit{embeddings} y \textit{reranking}) no está disponible, el flujo
principal debe seguir funcionando. En ese modo la capa semántica no participa:
no hay recuperación semántica, \textit{reranking} ni escritura indexada de
memorias de simulación en Qdrant.
\item \textbf{RNF-06: Determinismo.} Para una misma semilla
y configuración, el motor de simulación debe producir resultados
idénticos, de modo que las comparativas entre modelos sean
reproducibles y comparables entre sesiones.
\item \textbf{RNF-07: Portabilidad de despliegue.} Los servicios externos (PostgreSQL,
MinIO, Neo4j, Qdrant) deben poder levantarse de forma reproducible
mediante \texttt{docker compose}, sin pasos manuales de configuración
específicos del entorno del desarrollador.
\item \textbf{RNF-08: Configuración por entorno.} Las credenciales y \textit{endpoints} de
los servicios externos deben aislarse en variables de entorno, sin
quedar embebidos en el código fuente ni en los artefactos versionados.
\item \textbf{RNF-09: Observabilidad.} Las acciones del Orchestrator y de los
agentes especializados deben dejar traza en los \textit{logs} del
\textit{backend} (emitidos por la salida estándar del proceso que sirve la
API), suficientes para diagnosticar una sesión \textit{a posteriori} sin
tener que reproducirla.
\item \textbf{RNF-10: Acotación del contexto conversacional.} Una sesión larga no
debe agotar la ventana de contexto del modelo. El historial del chat con
el Orchestrator debe mantenerse dentro de un presupuesto acotado de forma
automática, sin que el usuario tenga que reiniciar la sesión y sin perder
la trazabilidad de las decisiones ya tomadas.
\item \textbf{RNF-11: Recuperación relacional segura.} Las consultas en
lenguaje natural sobre experimentos, modelos, observaciones y mensajes
persistidos deben traducirse a SQL de solo lectura, limitarse a un conjunto
explícito de tablas permitidas y aplicar un límite máximo de resultados. Si
la pregunta queda fuera de ese alcance, el sistema debe rechazarla sin
ejecutar SQL.
\item \textbf{RNF-12: Contraste previo con predicciones.} Antes de ejecutar una
simulación, el sistema debe poder recuperar las predicciones científicas
asociadas al modelo o paradigma seleccionado en la primera fase, de modo que el
análisis posterior pueda comparar la trayectoria observada con una expectativa
formulada antes de ver los resultados.
\end{itemize}
\end{small}

\clearpage
\section{Trazabilidad de requisitos}
\label{ap:trazabilidad-requisitos}

La tabla siguiente conecta los objetivos operativos de esta segunda fase
con los requisitos anteriores y con las evidencias usadas después en el
capítulo de pruebas. Su finalidad es evitar que los requisitos queden como
una lista declarativa: cada objetivo se vincula con una forma concreta de
comprobar que el sistema lo cumple.

\begin{table}[H]
\begin{center}
\small
\setlength{\tabcolsep}{5pt}
\renewcommand{\arraystretch}{1.15}
\begin{tabular}{|>{\raggedright\arraybackslash}p{3.6cm}|>{\raggedright\arraybackslash}p{2.7cm}|>{\raggedright\arraybackslash}p{5.7cm}|} \hline
\textbf{Objetivo} & \textbf{Requisitos} & \textbf{Evidencia} \\ \hline
Cargar modelos de la primera fase sin acoplamiento rígido & RF-01, RNF-01, RNF-02 & Pruebas del \texttt{model\_loader}, contrato \texttt{DecisionModel} y escenario multi-modelo. \\ \hline
Generar y validar entornos de simulación & RF-02, RF-03, RNF-06 & Tests de \texttt{simlab.spec} y \texttt{simlab.environment}; escenario con rejilla \(8\times8\), semilla fija y 30 pasos. \\ \hline
Observar trayectorias y eventos de decisión & RF-04, RF-12, RNF-03 & Salida estructurada del Tracker, trazas de decisión por agente y pruebas de emisión de eventos en vivo mediante WebSocket. \\ \hline
Analizar resultados y contrastarlos con conocimiento previo & RF-05, RF-06, RF-08, RNF-12 & \texttt{retrieve\_context}, pre-recuperación del Orchestrator, \texttt{read\_predictions} y gráficas del Analyst. \\ \hline
Generar informes revisables & RF-07 & Informe PDF del Reporter, plantilla LaTeX, descarga REST y pruebas de generación y descarga de informes. \\ \hline
Persistir conversaciones, experimentos y memorias de simulación & RF-08, RF-09, RF-10, RNF-03, RNF-10, RNF-11 & Tablas persistidas, \texttt{query\_history}/NL-SQL, autocompactación y escritura \(7=7=7\) en PostgreSQL y Qdrant. \\ \hline
Ofrecer una interfaz reactiva de uso & RF-10, RF-11, RF-12, RF-13, RNF-04 & Playwright, vistas de inspección, capturas de interfaz y verificación manual sobre la aplicación viva. \\ \hline
\end{tabular}
\caption[Trazabilidad de requisitos]{Trazabilidad de requisitos. Cada objetivo
operativo se vincula con sus requisitos y evidencias de validación.}
\end{center}
\end{table}
